\section{Chebyshev improvement from a common polynomial family}

Let $\mathcal H_M$ and $\mathcal H_E$ be finite-dimensional complex inner
product spaces, and let $H$ be self-adjoint on
$\mathcal H=\mathcal H_M\otimes\mathcal H_E$. Choose an orthonormal eigenbasis
$b_j$ with eigenvalues $h_j$. Suppose that $\eta_0<\eta_1\le u$, that
$h_j\le u$, and that each $h_j$ lies in $(-\infty,\eta_0]\cup[\eta_1,\infty)$.
For a nonnegative integer $k$, set
\begin{align}
    G&=\operatorname{span}\{b_j:h_j\le\eta_0\},
    \label{eq:cheb_viable_low_space}\\
    \Delta&=\bigl(2e^{-2k\sqrt{(\eta_1-\eta_0)/(u-\eta_0)}}\bigr)^2.
    \label{eq:cheb_viable_error}
\end{align}
These are the closed spectral conditions used in error reduction; positivity
of the unsquared filter is not required
\cite{gap:alvv16_spectral_projector_boundary,gap:alvv16_chebyshev_filter_positivity}.

\begin{theorem}[Chebyshev improvement from a uniform polynomial reduction]
    \label{thm:chebyshev_viability_polynomial_reduction}
    \lean{Submodule.chebyshev_viability_of_polynomial_reduction}
    \uses{thm:chebyshev_spectral_projector, thm:spectral_projector_contract}
    Fix a nonnegative integer $R$, a map $\Gamma$ on subspaces of
    $\mathcal H_M$, and an operator-valued map $\mathcal F$ on complex
    polynomials such that $\mathcal F(Q)=Q(H)$ for every $Q$.
    Suppose the following reduction holds for every $Q$ of degree at most
    $k$: whenever $S$ is $\delta$-viable for a subspace $T$ of $\mathcal H$,
    $\delta<1$, and $d\ge0$, and $\mathcal F(Q)$ preserves $T$ and $T^\perp$
    and satisfies, for every $t\in T$ and $z\in T^\perp$,
    \begin{align}
        \|\mathcal F(Q)t\|&\ge\|t\|,\\
        \|\mathcal F(Q)z\|^2&\le d\|z\|^2,
    \end{align}
    the subspace $\Gamma(S)$ has dimension at most $R\dim S$ and is
    $d/(1-\delta)^2$-viable for $T$.
    Then, under the spectral assumptions above, for every $\delta$-viable
    subspace $S$ for $G$ with $\delta<1$, the subspace $\Gamma(S)$
    has dimension at most $R\dim S$ and is
    $\Delta/(1-\delta)^2$-viable for $G$.
    Here $\Gamma$ and $\mathcal F$ are fixed before the spectral endpoints
    and before $S$. This conditional composition uses
    Theorem~4.2 of~\cite[lines 789--853]{Arad2016RigorousRG}
    and the error-reduction argument of Lemma~2.9.
\end{theorem}

\begin{proof}
    Choose the normalized real Chebyshev polynomial $p$ of degree $k$ and
    regard it as a complex polynomial. Its evaluation is a spectral
    projector with parameter~(\ref{eq:cheb_viable_error}). The closed spectral
    gap gives invariance, expansion on $G$, and contraction on $G^\perp$.
    Apply the supplied reduction to this polynomial and these bounds.
\end{proof}

\begin{theorem}[A common matrix family gives Chebyshev improvement]
    \label{thm:matrix_family_chebyshev_viability}
    \lean{Matrix.chebyshev_viability_of_conjugated_polynomial_family}
    \uses{thm:viable_tensor_decomposition, thm:chebyshev_viability_polynomial_reduction}
    Let $\Sigma,X,Y$ be finite sets, put $\mathcal H_M=\mathbb C^X$ and
    $\mathcal H_E=\mathbb C^Y$, and fix a unitary linear isometry
    $T:\mathbb C^\Sigma\longrightarrow\mathcal H_M\otimes\mathcal H_E$.
    Fix a matrix $H_0$ on $\mathbb C^\Sigma$ and matrices
    $W_1,\ldots,W_N$ on $\mathbb C^X$. Suppose that for every complex
    polynomial $Q$ of degree at most $k$ there are matrices
    $M_1(Q),\ldots,M_N(Q)$ on $\mathbb C^Y$ such that
    \begin{align}
        TQ(H_0)T^{-1}&=\sum_{r=1}^N W_r\otimes M_r(Q).
        \label{eq:cheb_viable_family}
    \end{align}
    Set $H=TH_0T^{-1}$ and $\Gamma(S)=\sum_{r=1}^N W_rS$.
    If $H$ satisfies the self-adjointness and spectral assumptions stated
    above, then every $\delta$-viable $S$ for
    $G$ of~(\ref{eq:cheb_viable_low_space}), with $\delta<1$, satisfies
    \begin{align}
        \dim\Gamma(S)&\le N\dim S,
        \label{eq:cheb_viable_dimension}\\
        \Gamma(S)&\text{ is }\Delta/(1-\delta)^2\text{-viable for }G.
        \label{eq:cheb_viable_conclusion}
    \end{align}
    The family $(W_r)$ is fixed before the spectral data and before $S$;
    only the complementary matrices $M_r(Q)$ depend on $Q$.
    The supplied polynomial identity is an explicit hypothesis. This is a
    finite-dimensional consequence supporting the polynomial decomposition
    and error-reduction steps of~\cite[Proposition 5.7 and Lemma 2.9]{Arad2016RigorousRG}.
\end{theorem}

\begin{proof}
    Polynomial evaluation commutes with the matrix action on Euclidean
    space and with conjugation by $T$. Thus the filter is $Q(H)$.
    For each $Q$, apply error reduction from its tensor decomposition
    in~(\ref{eq:cheb_viable_family}). The uniform polynomial reduction
    then gives~(\ref{eq:cheb_viable_dimension}) and~(\ref{eq:cheb_viable_conclusion}).
\end{proof}

\begin{theorem}[A numerical bound for the image sum]
    \label{thm:matrix_family_chebyshev_viability_card_bound}
    \lean{Matrix.chebyshev_viability_of_conjugated_polynomial_family_of_card_le}
    \uses{thm:matrix_family_chebyshev_viability}
    Under the hypotheses of the common matrix family theorem, let $C$ be a
    nonnegative integer with $N\le C$. The same fixed subspace
    $\Gamma(S)$ then satisfies $\dim\Gamma(S)\le C\dim S$ and is
    $\Delta/(1-\delta)^2$-viable for $G$, for every $\delta$-viable $S$
    with $\delta<1$.
    The numerical bound $C$ is fixed with the family before the spectral
    endpoints and before $S$.
\end{theorem}

\begin{proof}
    The dimension estimate follows from
    $\dim\Gamma(S)\le N\dim S\le C\dim S$. The viable space and its error
    are unchanged.
\end{proof}
